\section{Практические следствия, предсказательная сила и применение OC Core~1.3}
\label{sec:oc-core-1-3-practical-utility}

OC Core~1.3 практически полезна тогда, когда читателю нужно выбрать законный
пакет, сделать ограниченное остаточное предсказание, провести скрининг
аномалии, проверить переход состояния, наблюдать смену режима или сравнить
маршруты между доменами. Эта глава фиксирует, что уже можно использовать,
какие входы нужны, какие выходы это даёт и где проходит граница. Она не
представляет OC как неограниченную универсальную систему предсказания.

Практическая ценность задаётся playbook-ами и сравнительными матрицами. Они
показывают, что прежняя наука уже делает хорошо, где OC добавляет выбор
маршрута, тесты фальсификаторов и дисциплину replay, а где atlas честно
помечает комплементарность, интеграцию или превосходство локальной модели в
своём ограниченном коридоре.

Подробные row-level anchors находятся в
Appendix~\ref{sec:oc-core-1-3-practical-utility-atlas}. Для экономного
переводческого прохода канонические generated tables ниже сохранены в
машинно-верной форме; перевод этой главы отвечает за рамку применения, а не
за изменение чисел, формул, verdicts или trace identifiers.

\begingroup
\small
% Generated from OC_CORE_1_3_SCIENCE_SPOT_latest.json
\subsection{Why OC is useful}
The practical utility atlas currently tracks \occode{6} usable-now rows, \occode{3} frontier-program rows, and \occode{1} hypothesis-only row.
OC currently exposes 6 bounded practical lanes that are already usable under explicit theorem, data-route, and falsifier discipline.
Its closed practical value is not unrestricted universal prediction; it is lawful packet selection, bounded residual prediction, anomaly screening, state-transition auditing, and regime-shift monitoring under one source-bound route.
The atlas also keeps 3 frontier-program rows and 1 hypothesis-only row explicit so practical ambition does not masquerade as finished closure.

\subsection{Support classes}
Every practical claim below is explicitly labeled so the reader can see what is already usable, what is bounded, and what still remains exploratory.
\begin{itemize}[leftmargin=1.6em,nosep]
\item \occode{FORMALLY_PROVED} means formally proved and is usable now within the declared bounds.
\item \occode{THEOREM_NATIVE_HELD_OUT_VALIDATED} means theorem-native + empirical / held-out validated and is usable now within the declared bounds.
\item \occode{OPERATIONALLY_SUPPORTED_WITHIN_BOUNDS} means operationally supported within bounds and is usable now within the declared bounds.
\item \occode{FRONTIER_PROGRAM} means frontier program and is not yet a lawful usable-now claim.
\item \occode{HYPOTHESIS_ONLY} means hypothesis only and is not yet a lawful usable-now claim.
\end{itemize}

\subsection{What is already usable now}
\begingroup
\footnotesize
\sloppy
\setlength{\tabcolsep}{3pt}
\setlength{\LTleft}{0pt}
\setlength{\LTright}{0pt}
\begin{longtable}{@{}L{0.13\textwidth}L{0.18\textwidth}L{0.22\textwidth}L{0.13\textwidth}L{0.22\textwidth}@{}}
\toprule
Domain & Problem class & What OC gives now & Support & Output / decision \\
\midrule
\endfirsthead
\toprule
Domain & Problem class & What OC gives now & Support & Output / decision \\
\midrule
\endhead
Mathematics & Strict closure certification and counterexample elimination & Certify whether a bounded formal object or question packet survives strict closure and counterexample search under the locked mathematical lane. & formally proved & A replayable accept or reject result together with minimal denominator, augmentation rank, or surviving counterexample classification. \\
Physics & Constants, spectral-line, and transport-envelope audit & Audit bounded constants, Balmer-line families, and transport-scaling envelopes on held-out benchmark families and flag out-of-envelope residual behavior. & theorem-native + empirical / held-out validated & A pass or fail residual verdict, anomaly prioritization, and a decision about whether the family remains inside the promoted packet. \\
Chemistry & Thermochemical, spectroscopic, and kinetic packet audit & Audit bounded thermochemical, spectral, and reaction-envelope families on held-out compounds or reactions and flag out-of-envelope chemistry lanes. & theorem-native + empirical / held-out validated & A pass or fail residual verdict, compound or reaction-family screening result, and a decision about whether the lane remains inside the promoted chemistry packet. \\
Biology & Expression-signature, state-transition, and response-onset screening & Screen bounded biological lanes for expression-signature residuals, state-transition ordering errors, and response-onset residuals across held-out datasets. & theorem-native + empirical / held-out validated & A pass or fail signature verdict, a state-transition ordering judgment, and an escalation decision for the biological lane. \\
Systems / Civilizational projection & Macro-trajectory and regime-shift monitoring & Monitor bounded macro trajectories, turning-point candidates, and regime-shift signatures on held-out official series and flag out-of-band intervals. & theorem-native + empirical / held-out validated & A pass or fail trajectory verdict, a turning-point or regime-shift alert, and a decision about whether the lane remains inside the promoted systems packet. \\
Cross-domain / unified science & Cross-domain route selection, escalation discipline, and packet coordination & Choose the lawful domain packet, comparator set, and escalation route for a mixed scientific or operational problem without letting rhetoric outrun the closed packet boundary. & operationally supported within bounds & A lawful packet-selection decision, a comparator shortlist, and an explicit escalation path when the closed packet is not enough. \\
\bottomrule
\end{longtable}
\endgroup

\subsection{How to use OC in practice}
The playbooks below answer the operational question directly: when to use OC, what inputs are required, what procedure to run, and what the resulting decision or prediction actually is.

\paragraph{Mathematics: Strict closure and counterexample playbook}
When to use: Use when the task is to determine whether a bounded exact question already lies inside the promoted mathematical anchor or should be rejected by the strict counterexample route.
Inputs: formal question encoding; \allowbreak strict closure policy; \allowbreak counterexample search route.
Procedure: Encode the target object in the admissible exact-question form. $\rightarrow$ Run the declared terminalization and counterexample search. $\rightarrow$ Accept the result only if the replay is deterministic and survivor-free. $\rightarrow$ Escalate to frontier if the question leaves the declared exact scope.
Output: A replayable proof result, a surviving counterexample, or a lawful frontier demotion.
Support class: formally proved.
Failure boundary: Stop if the object cannot be encoded in the admissible exact-question format or if a surviving counterexample remains after the declared search.
Appendix~R carries the exact trace anchors and the full benchmark alignment for this playbook.

\paragraph{Physics: Bounded constants, spectra, and transport audit}
When to use: Use when the task is to decide whether a constants family, spectral family, or transport family belongs inside the promoted physics packet or should be escalated as an anomaly.
Inputs: chosen benchmark family; \allowbreak official dataset route; \allowbreak held-out split and sigma policy.
Procedure: Map the observable family to the promoted physics packet. $\rightarrow$ Load the pinned official dataset and the locked held-out split. $\rightarrow$ Run the declared residual replay and inspect sign or order constraints. $\rightarrow$ Accept the family only if the replay stays inside the promoted tolerance band.
Output: A bounded inclusion, anomaly, or escalation verdict for the target physical family.
Support class: theorem-native + empirical / held-out validated.
Failure boundary: Stop and demote to frontier if the observable family requires a widened scope, hidden fit freedom, or a post-hoc split change.
Appendix~R carries the exact trace anchors and the full benchmark alignment for this playbook.

\paragraph{Chemistry: Thermochemical, spectral, and kinetic packet audit}
When to use: Use when the task is to decide whether a bounded compound family or reaction class lies inside the promoted chemistry packet or should be escalated as an out-of-envelope case.
Inputs: compound or reaction-family definition; \allowbreak pinned NIST or DOE-linked route; \allowbreak held-out rotation and tolerance policy.
Procedure: Map the chemistry lane to the promoted packet and observable family. $\rightarrow$ Load the pinned reference tables and the locked held-out split. $\rightarrow$ Run the declared replay on thermochemical, spectral, or kinetic outputs. $\rightarrow$ Accept the lane only if the held-out family remains inside the promoted tolerance band.
Output: A bounded inclusion, anomaly, or escalation verdict for the target chemistry lane.
Support class: theorem-native + empirical / held-out validated.
Failure boundary: Stop and demote to frontier if the lane requires hidden fit freedom, widened scope, or post-hoc split changes.
Appendix~R carries the exact trace anchors and the full benchmark alignment for this playbook.

\paragraph{Biology: Bounded biological signature and onset screening}
When to use: Use when the task is to decide whether a bounded assay lane belongs inside the promoted biological packet or should be escalated as an unresolved biological signature family.
Inputs: bounded assay family; \allowbreak multi-dataset held-out split; \allowbreak declared signature, transition, or onset observable.
Procedure: Map the biological lane to the promoted packet and observable family. $\rightarrow$ Load the locked multi-dataset route and held-out split. $\rightarrow$ Run the declared replay on expression, transition-order, or onset outputs. $\rightarrow$ Accept the lane only if the held-out datasets stay inside the promoted signature envelope.
Output: A bounded inclusion, anomaly, or escalation verdict for the target biological lane.
Support class: theorem-native + empirical / held-out validated.
Failure boundary: Stop and demote to frontier if the lane depends on a single dataset family, hidden fit freedom, or a widened biological claim boundary.
Appendix~R carries the exact trace anchors and the full benchmark alignment for this playbook.

\paragraph{Systems / Civilizational projection: Bounded macro-trajectory and regime-shift monitoring}
When to use: Use when the task is to decide whether an official macro-process series remains inside the promoted systems packet or should be escalated as a turning-point or regime anomaly.
Inputs: target macro-process series; \allowbreak held-out interval; \allowbreak declared trajectory or turning-point metric.
Procedure: Map the macro-process lane to the promoted systems packet. $\rightarrow$ Load the pinned official series and the locked held-out interval. $\rightarrow$ Run the declared residual and turning-point replay. $\rightarrow$ Accept the lane only if held-out intervals remain inside the promoted trajectory and regime envelope.
Output: A bounded inclusion, anomaly, or escalation verdict for the target systems lane.
Support class: theorem-native + empirical / held-out validated.
Failure boundary: Stop and demote to frontier if the lane requires out-of-band trajectory behavior, a widened horizon, or post-hoc interval selection.
Appendix~R carries the exact trace anchors and the full benchmark alignment for this playbook.

\paragraph{Cross-domain / unified science: Cross-domain packet selection and escalation playbook}
When to use: Use when a scientific or engineering problem touches more than one domain lane and the team must decide which closed packet, comparator family, and evidence route should govern the next action.
Inputs: declared target problem class; \allowbreak bounded observable family; \allowbreak candidate domain packets; \allowbreak decision or intervention objective.
Procedure: Map the target problem to the narrowest lawful observable family. $\rightarrow$ Select the closed packet whose theorem-to-observable map already governs that family. $\rightarrow$ Check same-claim baselines and serious comparators before claiming advantage. $\rightarrow$ Escalate only when the closed packet boundary is reached and the atlas permits the translation.
Output: A lawful route-selection decision with explicit packet, comparator, falsifier, and escalation boundaries.
Support class: operationally supported within bounds.
Failure boundary: Stop and demote to frontier if the observable family is outside every closed packet or if two candidate routes produce incompatible bounded verdicts.
Appendix~R carries the exact trace anchors and the full benchmark alignment for this playbook.

\subsection{What OC predicts by domain}

\paragraph{Physics.}
The closed packet currently predicts or audits the following bounded lane(s): Audit bounded constants, Balmer-line families, and transport-scaling envelopes on held-out benchmark families and flag out-of-envelope residual behavior.
Measured outputs are \occode{relative_constant_error}, \occode{spectral_line_residual}, \occode{transport_scaling_residual}; benchmark families are PHY\_BENCH\_001; fundamental constants residuals; spectral line residual envelopes; transport or plasma scaling envelopes; held-out case total is 20.
What remains open: The repository does not yet lawfully promote a general OC replacement for the full physics model stack.

\paragraph{Chemistry.}
The closed packet currently predicts or audits the following bounded lane(s): Audit bounded thermochemical, spectral, and reaction-envelope families on held-out compounds or reactions and flag out-of-envelope chemistry lanes.
Measured outputs are \occode{enthalpy_residual}, \occode{spectral_transition_residual}, \occode{kinetic_or_equilibrium_residual}; benchmark families are CHEM\_BENCH\_001; thermochemical reference values; spectroscopic transition families; kinetic or equilibrium observable envelopes; held-out case total is 19.
What remains open: The repository does not yet promote a general OC catalyst-discovery or unrestricted synthetic-chemistry engine.

\paragraph{Biology.}
The closed packet currently predicts or audits the following bounded lane(s): Screen bounded biological lanes for expression-signature residuals, state-transition ordering errors, and response-onset residuals across held-out datasets.
Measured outputs are \occode{expression_signature_residual}, \occode{state_transition_order_error}, \occode{response_onset_residual}; benchmark families are BIO\_BENCH\_001; gene-expression time-series families; cell-state transition signatures; stimulus-response or excitability onset markers; held-out case total is 16.
What remains open: The repository does not yet promote a general OC biological intervention or whole-phenotype control engine.

\paragraph{Systems / Civilizational projection.}
The closed packet currently predicts or audits the following bounded lane(s): Monitor bounded macro trajectories, turning-point candidates, and regime-shift signatures on held-out official series and flag out-of-band intervals.
Measured outputs are \occode{trajectory_residual}, \occode{turning_point_error}, \occode{held_out_interval_regime_score}; benchmark families are SYS\_BENCH\_001; macro indicator trajectories; resource-intensity or throughput curves; held-out interval stability and regime shifts; held-out case total is 18.
What remains open: The repository does not yet promote unrestricted long-horizon civilizational forecasting as lawful closed science.

\subsection{Benchmark against serious alternatives}
This chapter distinguishes between simpler same-claim-class baselines and serious established model families. The first question is whether a cheaper competitor already solves the same bounded lane. The second is whether OC merely coexists with, integrates, or improves the fragmented scientific picture around that lane.

\paragraph{Same-claim-class baselines.}
\begingroup
\footnotesize
\sloppy
\setlength{\tabcolsep}{3pt}
\setlength{\LTleft}{0pt}
\setlength{\LTright}{0pt}
\begin{longtable}{@{}L{0.14\textwidth}L{0.24\textwidth}L{0.13\textwidth}L{0.37\textwidth}@{}}
\toprule
Domain & Comparator families & Verdict & Comparison scope \\
\midrule
\endfirsthead
\toprule
Domain & Comparator families & Verdict & Comparison scope \\
\midrule
\endhead
Mathematics & exact enumerative baseline; direct constructive proof baseline & \occode{DECLARED_NO_SUPERIOR_BASELINE} & Theorem-native anchor lane is trace-closed and replayable. \\
Physics & bounded residual fit without theorem trace; family-wise least-squares baseline & \occode{NO_SIMPLER_SUPERIOR_BASELINE_FOUND} & Bounded theorem-native physics packet is trace-closed, replay-closed, and comparator-clean on the pinned official held-out route. \\
Chemistry & reference-table envelope baseline; reaction-class residual fit baseline & \occode{NO_SIMPLER_SUPERIOR_BASELINE_FOUND} & Bounded theorem-native chemistry packet is trace-closed, replay-closed, and comparator-clean on the pinned thermochemical, spectral, and kinetic held-out route. \\
Biology & dataset-family signature fit baseline; state-transition score baseline & \occode{NO_SIMPLER_SUPERIOR_BASELINE_FOUND} & Bounded theorem-native biology packet is trace-closed, replay-closed, and comparator-clean across the pinned multi-dataset held-out biological lanes. \\
Systems / Civilizational projection & macro-trend residual baseline; turning-point heuristic baseline & \occode{NO_SIMPLER_SUPERIOR_BASELINE_FOUND} & Bounded theorem-native systems packet is trace-closed, replay-closed, and comparator-clean on the pinned macro-trajectory and regime-shift held-out route. \\
\bottomrule
\end{longtable}
\endgroup

\paragraph{Serious model families.}
\begingroup
\footnotesize
\sloppy
\setlength{\tabcolsep}{3pt}
\setlength{\LTleft}{0pt}
\setlength{\LTright}{0pt}
\begin{longtable}{@{}L{0.13\textwidth}L{0.18\textwidth}L{0.15\textwidth}L{0.13\textwidth}L{0.33\textwidth}@{}}
\toprule
Domain & Model family & Problem class & Relation to OC & What OC adds or closes \\
\midrule
\endfirsthead
\toprule
Domain & Model family & Problem class & Relation to OC & What OC adds or closes \\
\midrule
\endhead
Physics & Gyrokinetic, confinement-scaling, and plasma-transport models & transport or plasma scaling envelopes & complements & OC adds a shared bounded packet and replay rule that lets transport residual families live beside constants and spectra without rhetorical widening. \\
Physics & Precision quantum electrodynamics and atomic-structure theory & fundamental constants and spectral transition families & integrates & OC adds one bounded theorem-to-observable packet that keeps constants, Balmer spectra, and transport residual families under one explicit closure, falsifier, and held-out replay contract. \\
Chemistry & Chemical kinetics, transition-state, and equilibrium models & reaction-class kinetic or equilibrium envelopes & integrates & OC adds a bounded chemistry packet that keeps those observable families under one theorem, one parameter-law shell, and one held-out replay discipline. \\
Chemistry & Quantum chemistry, ab initio, and density-functional families & thermochemical and spectroscopic envelopes & still superior in lane & OC adds a bounded packet that ties thermochemical, spectroscopic, and kinetic envelopes to one theorem-native route with one falsifier and one replay discipline. \\
Biology & Excitable-media and response-onset dynamical models & response-onset residual lanes & complements & OC adds a single bounded biological packet that keeps response onset tied to expression and transition-order evidence inside one replay and falsifier framework. \\
Biology & Gene-regulatory, transcriptomic, and systems-biology state-space models & expression-signature and state-transition lanes & integrates & OC adds a bounded packet that keeps expression signatures, state-transition order, and response onset under one multi-dataset held-out closure route. \\
Systems / Civilizational projection & Systems dynamics, control, and change-point models & macro-process trajectories and regime shifts & integrates & OC adds a bounded macro-trajectory packet that keeps trajectory residuals, turning points, and regime scores under one theorem-to-observable route and one held-out decision bar. \\
Cross-domain / unified science & Complexity-science and multiscale integration programs & cross-domain integration and multiscale explanation & complements & OC adds a harder closure grammar: explicit theorem packets, parameter laws, held-out replay, falsifiers, and same-claim comparator ledgers for each bounded lane. \\
Cross-domain / unified science & Reductionist patchwork of separate domain theories & cross-domain scientific planning and explanation & integrates & OC adds one explicit kernel, one K-level ladder, one theorem-to-observable grammar, and one atlas for translating bounded packets across domains without rhetorical leakage. \\
\bottomrule
\end{longtable}
\endgroup

\subsection{What prior science could do, what remained fragmented, and what OC closes}
\begingroup
\footnotesize
\sloppy
\setlength{\tabcolsep}{3pt}
\setlength{\LTleft}{0pt}
\setlength{\LTright}{0pt}
\begin{longtable}{@{}L{0.12\textwidth}L{0.20\textwidth}L{0.20\textwidth}L{0.23\textwidth}L{0.17\textwidth}@{}}
\toprule
Domain & Prior science already did & Fragmentation or limit & What OC closes or adds & Open boundary \\
\midrule
\endfirsthead
\toprule
Domain & Prior science already did & Fragmentation or limit & What OC closes or adds & Open boundary \\
\midrule
\endhead
Physics & Provides strong local transport and confinement heuristics for specific plasma regimes and benchmark devices. & These models are powerful in their native lane but do not by themselves close a unified residual grammar spanning constants, spectra, and transport under one promotion rule. & OC adds a shared bounded packet and replay rule that lets transport residual families live beside constants and spectra without rhetorical widening. & OC does not replace specialized transport solvers or device-specific plasma modeling outside the promoted packet. \\
Physics & Delivers high-precision local calculations for constants relations, atomic transitions, and spectral families inside explicit microscopic equation classes. & Selection and justification of the local model family remain lane-specific, and the route does not by itself give a shared kernel-to-observable closure discipline across domains. & OC adds one bounded theorem-to-observable packet that keeps constants, Balmer spectra, and transport residual families under one explicit closure, falsifier, and held-out replay contract. & Detailed local derivation outside the promoted residual packet still relies on the standard physics model stack. \\
Chemistry & Explains rate, barrier, and equilibrium behavior for declared reaction classes with strong domain-specific local formalisms. & The family remains reaction-class specific and does not itself unify thermochemistry, spectroscopy, and kinetics under one kernel-bound promotion route. & OC adds a bounded chemistry packet that keeps those observable families under one theorem, one parameter-law shell, and one held-out replay discipline. & OC does not erase the need for local kinetic modeling outside the promoted bounded families. \\
Chemistry & Provides molecule-specific energetic and spectroscopic calculations with a mature tool chain for local chemical systems. & Method choice, basis dependence, and domain-specific tuning remain lane-specific, and the family does not by itself provide a closure ladder from kernel operators to held-out benchmark governance. & OC adds a bounded packet that ties thermochemical, spectroscopic, and kinetic envelopes to one theorem-native route with one falsifier and one replay discipline. & High-precision molecule-specific calculation remains the job of the established chemistry stack outside the promoted OC packet. \\
Biology & Captures threshold and onset dynamics well in specific excitability, signaling, or response families. & The family remains local to its chosen mechanism and does not by itself close broader biological signature, transition, and onset lanes under one bounded acceptance bar. & OC adds a single bounded biological packet that keeps response onset tied to expression and transition-order evidence inside one replay and falsifier framework. & Detailed mechanism-specific onset modeling still belongs to the established lane-specific models outside the promoted packet. \\
Biology & Provides established state-space, regulatory-network, and transcriptomic modeling inside specific assay and dataset families. & Performance is often lane- and dataset-specific, and cross-dataset closure plus theorem-to-observable promotion discipline remain fragmented. & OC adds a bounded packet that keeps expression signatures, state-transition order, and response onset under one multi-dataset held-out closure route. & OC does not yet replace detailed mechanistic biological modeling outside the promoted assay lanes. \\
Systems / Civilizational projection & Provides strong local tools for trend extraction, control analysis, and regime-change detection in declared systems lanes. & These tools are typically lane-specific and do not by themselves close a shared theorem route from kernel operators to official macro-process held-out promotion. & OC adds a bounded macro-trajectory packet that keeps trajectory residuals, turning points, and regime scores under one theorem-to-observable route and one held-out decision bar. & OC does not promote unrestricted civilizational prediction or replace every local macro or control model outside the bounded packet. \\
Cross-domain / unified science & Provides rich ideas about emergence, multiscale organization, and complex adaptive behavior across many scientific domains. & These programs are often strong conceptually but heterogeneous in proof bar, operator discipline, and closed promotion criteria. & OC adds a harder closure grammar: explicit theorem packets, parameter laws, held-out replay, falsifiers, and same-claim comparator ledgers for each bounded lane. & OC does not exhaust every possible complexity-science question; it narrows attention to the lanes it can lawfully close. \\
Cross-domain / unified science & Explains many local scientific lanes well by using separate best-in-class models inside each domain. & Cross-domain translation, shared operator meaning, and unified promotion discipline remain fragmented across disconnected model families. & OC adds one explicit kernel, one K-level ladder, one theorem-to-observable grammar, and one atlas for translating bounded packets across domains without rhetorical leakage. & OC does not abolish local domain theories; it governs how bounded closed packets are related and promoted together. \\
\bottomrule
\end{longtable}
\endgroup

\subsection{What remains frontier or hypothesis}
\begingroup
\footnotesize
\sloppy
\setlength{\tabcolsep}{3pt}
\setlength{\LTleft}{0pt}
\setlength{\LTright}{0pt}
\begin{longtable}{@{}L{0.11\textwidth}L{0.19\textwidth}L{0.14\textwidth}L{0.48\textwidth}@{}}
\toprule
Domain & Problem class & Support & Why this is not yet closed \\
\midrule
\endfirsthead
\toprule
Domain & Problem class & Support & Why this is not yet closed \\
\midrule
\endhead
Physics & Broader microphysical law extension beyond the promoted residual packet & frontier program & The repository does not yet lawfully promote a general OC replacement for the full physics model stack. \\
Chemistry & General catalyst or out-of-family chemistry design & frontier program & The repository does not yet promote a general OC catalyst-discovery or unrestricted synthetic-chemistry engine. \\
Biology & Multi-scale phenotype control and unrestricted biological intervention & frontier program & The repository does not yet promote a general OC biological intervention or whole-phenotype control engine. \\
Systems / Civilizational projection & Unrestricted long-horizon civilizational forecasting and control & hypothesis only & The repository does not yet promote unrestricted long-horizon civilizational forecasting as lawful closed science. \\
\bottomrule
\end{longtable}
\endgroup

\endgroup

\paragraph{Переход к atlas.}
Практическая глава намеренно ограничена: она несёт high-signal таблицы для
использования читателем, а полный audit burden оставляет
Appendix~\ref{sec:oc-core-1-3-practical-utility-atlas}.
